
%-------------------------------------------------------------------------------
% FORMATAÇÃO
%-------------------------------------------------------------------------------

% Pacotes utilizados:
\usepackage{amsmath}
\usepackage{graphicx}
\usepackage[colorlinks=true, allcolors=blue]{hyperref}
\usepackage{hyperref}
\usepackage{orcidlink}
\usepackage[title]{appendix}
\usepackage{mathrsfs}
\usepackage{amsfonts}
\usepackage{booktabs}                                      % \toprule, \midrule, \botrule;
\usepackage{caption}                                       % \caption;
\usepackage{threeparttable}                                % inserir rodapé nas tabelas;
\usepackage{algorithm}
\usepackage{algorithmicx}
\usepackage{indentfirst}
\usepackage{algpseudocode}
\usepackage{listings}
\usepackage{enumitem}
\usepackage{chngcntr}
\usepackage{booktabs}
\usepackage{lipsum}
\usepackage{subcaption}
\usepackage{authblk}
\usepackage[T1]{fontenc}                                   % font encoding;
\usepackage{csquotes}                                      % csquotes;
\usepackage{diagbox}
\usepackage[
  style=abnt,
  backref=true,
  backend=biber,
  citecounter=true,
  backrefstyle=three, 
  url=true,
  maxbibnames=99,
  mincitenames=1,
  maxcitenames=3,
  backref=false,
  hyperref=true,
  giveninits=true,
  uniquename=false,
  uniquelist=false]{biblatex}
\usepackage[brazilian]{babel}                             % setar idioma.

% Margens:  
\usepackage[
  top=2cm,
  bottom=2cm,
  left=2.5cm,
  right=2.5cm,
  marginparwidth=1.75cm]{geometry}

% Referências:
\addbibresource{arquivos/referencias.bib}

% Mudar 'and' para '&' nas referências:
\renewcommand*{\finalnamedelim}{
  \ifnumgreater{\value{liststop}}{2}{\finalandcomma}{}
  \addspace\&\space}


%% Configurações:

% Espaçamento das linhas:
\usepackage{setspace}
\onehalfspacing

% Redefinir formato da numeração da seção e subseção:
\usepackage{titlesec}
\titleformat{\section}
  {\normalfont\Large\bfseries}{\thesection.}{1em}{}
  
% Customizar numeração das linhas:
\usepackage{lineno}
\renewcommand\linenumberfont{\normalfont\scriptsize\sffamily\color{blue}}
\rightlinenumbers
%\linenumbers                      % inserir numeração das linhas.

% Define um novo comando para fourth-level title.
\newcommand{\subsubsubsection}[1]{
  \vspace{\baselineskip}
  \noindent\textbf{#1\\}\quad
}

% Alterar a posição da legenda da tabela acima da tabela:
\usepackage{float}
\usepackage{caption}

% Definir lista não numerada:
\makeatletter
\newenvironment{unlist}{
  \begin{list}{}{
    \setlength{\labelwidth}{0pt}
    \setlength{\labelsep}{0pt}
    \setlength{\leftmargin}{2em}
    \setlength{\itemindent}{-2em}
    \setlength{\topsep}{\medskipamount}
    \setlength{\itemsep}{3pt}
  }
}{
  \end{list}
}
\makeatother

% Apêndice:
\renewcommand\theequation{\Alph{section}\arabic{equation}} % Redefine equation numbering format
\counterwithin*{equation}{section}                         % Number equations within sections
\renewcommand\thefigure{\Alph{section}\arabic{figure}}     % Redefine equation numbering format
\counterwithin*{figure}{section}                           % Number equations within sections
\renewcommand\thetable{\Alph{section}\arabic{table}}       % Redefine equation numbering format
\counterwithin*{table}{section}                            % Number equations within sections


